\documentclass[../main]{subfiles}
\begin{document}
\ifSubfilesClassLoaded{\setcounter{page}{54}}{}
\section{Комунізм і суспільна свідомість}
\label{sec:10_kommunizm_soznanie}

З марксистської точки зору, здійснюючи матеріальне виробництво, суспільство створює власну свідомість. Людство знає різні історично сформовані форми суспільної свідомості. Від міфологічної свідомості, де ідеальне ще не відділене або на різних стадіях не повністю відділене від безпосередньої матеріальної діяльності, з розвитком поділу праці та відчуження людини від власної діяльності виокремлюються філософія та релігія, потім мистецтво і наука. Паралельно з ними існують мораль і право, які виникли тоді, коли суспільство почало потребувати зовнішнього регулювання взаємовідносин між людьми з протилежними інтересами. Суспільна свідомість – це не сума свідомостей окремих індивідів, а свідомість суспільства як цілого, що відокремилася від безпосередньої людської діяльності, її ідеальний вимір. І навпаки, індивідуальна свідомість – це та сама суспільна свідомість, яку людина набуває, залежно від її місця та ролі в суспільному житті.

Суспільна свідомість та її форми мають об’єктивне існування щодо членів суспільства. Щоб жити в людському суспільстві, з його вимогами та настановами щодо моралі, права, релігії, науки, а також мистецтва та філософії, люди змушені рахуватися з ними так само, як і з навколишнім середовищем та рівнем розвитку технологій. У суспільній свідомості відображено і виражено не щось довільне, а спосіб взаємодії людини з природою та людини з людиною. Суспільна свідомість і всі її форми – не безпосередній продукт природи, а продукт її освоєння людським суспільством, відображений спосіб виробництва.

Те, що суспільна свідомість відокремилася від безпосередньої діяльності та набула відносно самостійного існування у своїх формах, а також те, що вона має обов’язковий характер для членів суспільства, було зафіксовано задовго до Маркса. Це лежить в основі не лише релігійного світогляду, а й усієї ідеалістичної філософії.

Формування ідеї комунізму вимагало критики як самого цього факту, так і того, як він інтерпретується в теоретичному мисленні. Критика суспільної свідомості не лише певної епохи, а й суспільства, де панує експлуатація людини людиною взагалі, відіграла найважливішу роль у формуванні ідейного змісту комунізму.

Як відомо, марксизм пройшов школу критики філософії Гегеля та фейєрбаховської критики християнської релігії – форми суспільної свідомості, яка була домінуючою в Європі протягом Середньовіччя і продовжує претендувати на домінування в буржуазному суспільстві. Наскільки серйозно Маркс ставився до фейєрбаховської критики релігії, видно з того, що він говорить про Фейєрбаха як про людину, яка здійснила неабиякий подвиг. В чому, на думку Маркса, полягає цей подвиг у зв’язку з нашим питанням? «1) у доведенні того, що філософія – це не що інше, як виражена в думках і логічно систематизована релігія, не що інше, як інша форма, інший спосіб існування відчуження людської сутності, і що, отже, вона також підлягає осуду (це – дуже важливе зауваження, що стосується філософії як форми суспільної свідомості); 2) в заснуванні істинного матеріалізму та реальної науки, оскільки суспільне відношення «людина до людини» Фейєрбах також робить основним принципом теорії».

Це переваги концепції Фейєрбаха, але вона не є повністю послідовною, і Маркс спеціально зазначив це: «Фейєрбах виходить із факту релігійного відчуження, з подвоєння світу на релігійний, уявний світ і дійсний світ. І він займається тим, що зводить релігійний світ до його земної основи. Він не помічає, що після виконання цієї роботи головне ще не зроблено. А саме, той факт, що земна основа відокремлюється від самої себе і переносить себе в хмари як певне самостійне царство, може бути пояснений лише саморозривом і самосуперечливістю цієї земної основи. Отже, остання, по-перше, повинна бути зрозуміла у своїй суперечності, а потім практично революціонізована шляхом усунення цієї суперечності. Отже, після того як, наприклад, у земній сім’ї знайдено розгадку таємниці святої сім’ї, сама земна сім’я повинна бути піддана теоретичній критиці та практично революційно перетворена. Отже, після того як, наприклад, у земній сім’ї знайдено розгадку таємниці святої сім’ї, сама земна сім’я має бути піддана теоретичній критиці та практично революційно перетворена.

Таким чином, земне життя людини – матеріальне виробництво – стає для Маркса тим предметом, який потрібно вивчати, щоб зрозуміти сутність земних коренів свідомості. «Релігія сама по собі позбавлена змісту, її витоки знаходяться не на небі, а на землі, і зі знищенням тієї збоченої реальності, теоретичним вираженням якої вона є, вона зникає сама собою». (Маркс – Арнольду Руге, Дрезден, 30 листопада 1842 р.)

Те саме стосується й інших форм суспільної свідомості: вже згаданої філософії, моралі, права, і навіть науки й мистецтва, їх потрібно розуміти як сторони людської практики, виходячи з неї самої, з її характеру та змісту, а не з них самих. У марксизмі чітко розрізняють базис суспільних відносин – єдність продуктивних сил і виробничих відносин – і похідну від них надбудову, до якої належить суспільна свідомість.

Наприклад, у підготовчих рукописах до «Капіталу» Маркс пише про право наступне: «Якщо звести ці тривіальності (уявлення про право сучасних Марксу економістів – ред.) до їхньої справжньої сутності, то вони скажуть більше, ніж знають ті, хто їх проповідує. А саме, що кожна форма виробництва породжує властиві їй правові відносини, форми правління тощо. У грубості та відсутності розуміння полягає те, що органічно пов’язані явища ставляться у випадкові взаємовідносини та у чисто логічний зв’язок. Буржуазним економістам здається, що при сучасній поліції можна краще виробляти, ніж, наприклад, при праві сильного. Вони просто забувають, що право сильного – це також право, і що воно в іншій формі продовжує існувати і в їхній «правовій державі».

«Земні корені» суспільної свідомості мають бути не лише зрозумілі, але й змінені – змінено характер суспільного виробництва. І те, що форми суспільної свідомості виростають із способу суспільного виробництва, свідчить про те, що з ними справи обстоюються зовсім не так просто, як може здатися на перший погляд. Наприклад, повністю скасувати право, подолати необхідність у ньому як зовнішньому регуляторі можна лише за певних суспільних умов – за комунізму, розвиненого на власній основі. Те саме стосується і подолання релігії. Її не можна просто скасувати, оголосивши її хибною. Її потрібно усунути, зберігши всі досягнення людської культури, які розвивалися в релігійній формі, знищивши шляхом докорінного зміни характеру щоденної людської діяльності.

Звісно, у невеликій брошурі про суспільну свідомість за комунізму ми можемо говорити лише в загальних рисах. Але навіть у цьому випадку слід відзначити одну найважливішу обставину, що стосується суспільства, яке існувало до комунізму і було поділене на класи. Так само, як люди займають різне становище в процесі виробництва, що відповідає їхній класовій приналежності, так і суспільна свідомість має класовий характер. Це не означає, що не існує єдиного поля свідомості та об’єктивних, незалежних виключно від класової позиції її носія, істин, добра і краси. Це означає, що саме це поле свідомості є полем класової боротьби. Це означає, що панівний економічний клас панує також і в свідомості. Що свідомість панівного класу є панівною в класовому суспільстві. Тому Ленін, слідуючи за Марксом і Енгельсом, говорить про важливість боротьби на полі свідомості – про теоретичну боротьбу, як про особливий вид класової боротьби.

У брошурі «Що робити?» Ленін присвятив цілий розділ значенню теоретичної боротьби пролетаріату та поглядам Енгельса з цього питання. На нашу думку, наступне зауваження Леніна є надзвичайно актуальним і сьогодні: «Усі, хто хоч трохи знайомий з фактичним станом нашого руху, не можуть не помітити, що широке поширення марксизму супроводжувалося певним зниженням теоретичного рівня. До руху, заради його практичного значення та практичних успіхів, приєднувалося чимало людей, які були дуже мало або взагалі не підготовлені теоретично. Тому можна судити про те, яке відсутність такту демонструє «Раб. Діло», коли з тріумфом представляє вислів Маркса: «кожен крок реального руху важливіший, ніж дюжина програм». Повторювати ці слова в епоху теоретичної розгубленості – це все одно, що кричати «вам нічим не допомогти!» при вигляді похоронної процесії. Та й ці слова Маркса взято з його листа про Готську програму, в якому він різко критикує допущену еклектику у формулюванні принципів: якщо вже треба було об’єднуватися, – писав Маркс лідерам партії, – то укладайте угоди заради задоволення практичних цілей руху, але не допускайте торгівлі принципами, не робіть теоретичних «компромісів». Ось яка була думка Маркса, а у нас є люди, які, в ім’я його, намагаються послабити значення теорії! Без революційної теорії не може бути і революційного руху. Не можна достатньо наголошувати на цій думці в такий час, коли з модною проповіддю оппортунізму поєднується захоплення найвужчими формами практичної діяльності.

Звертаючись до спадщини Енгельса, Ленін наводить велику цитату, яку ми вважаємо за потрібне також навести і тут: «Німецькі робітники мають дві суттєві переваги перед робітниками решти Європи. Перша полягає в тому, що вони належать до найбільш теоретичного народу Європи і що вони зберегли в собі той теоретичний зміст, який майже повністю втрачено так званими «освіченими» класами в Німеччині. Без попередньої німецької філософії, особливо філософії Гегеля, ніколи б не було створено німецький науковий соціалізм, – єдиний науковий соціалізм, який коли-небудь існував. Без теоретичного змісту у робітників цей науковий соціалізм ніколи б не проник у їхню плоть і кров у такій мірі, як ми бачимо зараз». А наскільки великою є ця перевага, це видно, з одного боку, з тієї байдужості до будь-якої теорії, яка є однією з головних причин того, чому англійський робітничий рух так повільно просувається вперед, незважаючи на чудову організацію окремих ремесел, — а з іншого боку, це показує ту плутанину та ті коливання, які посіяв прудонізм у своїй початковій формі у французів і бельгійців, а також у своїй карикатурній формі, наданій Бакуніним, — в Іспанії та Італії.

Друга перевага полягає в тому, що німці долучилися до робітничого руху майже в останню чергу. Як німецький теоретичний соціалізм ніколи не забуде, що він стоїть на плечах Сен-Сімона, Фурьє та Оуена — трьох мислителів, які, попри всю фантастичність і утопізм їхніх вчень, належать до найвидатніших умів усіх часів і які геніально передбачили незліченну кількість істин, правильність яких ми тепер доводимо науково, — так само німецький практичний робітничий рух ніколи не повинен забувати, що він розвинувся на плечах англійського та французького руху, що він мав можливість просто використати на свою користь їхній дорогою здобутий досвід, уникнути тепер їхніх помилок, яких тоді в більшості випадків неможливо було уникнути. Де б ми були зараз без прикладу англійських профспілок і французької політичної боротьби робітників, без того колосального поштовху, який особливо дала Паризька Комуна?

Слід віддати належне німецьким робітникам за те, що вони з надзвичайною майстерністю використали можливості, які випливали з їхнього становища. Вперше з часу існування робітничого руху боротьба ведеться систематично за всіма трьома її напрямками, узгодженими та взаємопов’язаними: в теоретичному, політичному та практично-економічному (опір капіталістам).

Це місце показує ставлення Енгельса до теоретичної боротьби. Але класова боротьба в царині суспільної свідомості не обмежується лише теоретичною боротьбою, так само як і суспільна свідомість не обмежується понятійним мисленням. Боротьба ведеться також у сфері естетичної діяльності, у сфері моралі, а також щодо права та релігії. Це також не можна недооцінювати. Однак, теоретична робота, спрямована на набуття та розвиток пролетаріатом власної класової свідомості, є ключем до всіх інших форм свідомості, оскільки дозволяє їх зрозуміти.

Варто також зазначити, що в класовому суспільстві в суспільній свідомості, незалежно від того, чи усвідомлюють люди це, загальнолюдське реалізується через класове. Свідомість революційного класу, який на даному історичному етапі виражає загальнолюдський інтерес, є рушійною силою історії та реалізує це загальнолюдське. У цьому сенсі не лише мораль, релігія та філософія, право, а й наука та мистецтво є також ідеологією, оскільки як естетичне освоєння світу в мистецтві, так і пізнання в науці, зумовлені, а отже й обмежені (свобода, яка в них присутня, обмежена) історичним розвитком суспільства та розстановкою сил у класовій боротьбі.

У зв’язку з цим хотілося б звернутися до питання про суб’єкта мислення у Маркса та Леніна. Слід нагадати, що Маркс поділяв політичну економію на класичну та вульгарну, використовуючи для цього цілком певний критерій. Він вказує на те, що класична політекономія була етапом розвитку розуміння свого предмета, на відміну від вульгарної. Він чітко визначає партійність, тобто класовість цієї науки. І класична, і вульгарна політична економія – це буржуазна наука, яка стоїть на службі у буржуазії та відображає в теорії напрямок її розвитку. Як тільки цей розвиток перестає збігатися з розвитком суспільства, перестає втілювати цей розвиток, іншими словами, коли буржуазія перестає бути революційним класом, політична економія стає вульгарною. Це означає, що вона вже не є такою, що розвивається, а навпаки, стає перешкодою на шляху розвитку думки.

Звісно, справа тут не в особистих якостях окремих політекономів. Маркс переконливо показує, що суб’єктом мислення є не окремі теоретики і не суспільство в цілому, а саме клас. Більше того, він є суб’єктом мислення лише до тих пір, поки він є суб’єктом історичної дії.

Ту саму думку ми знаходимо і в Леніна, коли він говорить про партії у філософії, яких може бути лише дві: партія, що мислить, і партія, яка перешкоджає мисленню. І саме тому, що одна партія – це партія класу, який є суб’єктом історичної дії, а отже, і суб’єктом мислення, а інша – ні. Це стосується не лише філософії, а ВСІХ без винятку наук, зокрема й природничих. Хоча там це не так очевидно, але, наприклад, отримати фінансування для досліджень, які не приносять навіть у перспективі жодних плодів, стає настільки важко, що це цілком можна розглядати як епізоди завуальованої класової боротьби. Це залишається актуальним до тих пір, поки розвиток суспільства здійснюється через класову боротьбу. Виходить, що об’єктивна істина завжди тенденційна, має класовий характер. Те саме стосується об’єктивності добра і краси.

У вересні (за старим стилем – у серпні) 1917 року Ленін висловив добре відому і, на нашу думку, дуже важливу думку, яка перегукується з вищесказаним і показує, як у революційні епохи пролетаріат ставиться до своєї класової свідомості, враховуючи, що мета пролетаріату, на відміну від інших класів, полягає не в тому, щоб завоювати і закріпити своє класове панування, а в тому, щоб, завоювавши це панування, знищити будь-яке класове панування і будь-який поділ суспільства на класи: «Будемо непохитні, – каже Ленін, – у розборі найменших сумнівів судом свідомих робітників, судом своєї партії, їй ми віримо, в ній ми бачимо розум, честь і совість нашої епохи, в міжнародному союзі революційних інтернаціоналістів ми бачимо єдину гарантію визвольного руху робітничого класу». «Панування класу неможливе без панування в сфері свідомості», – наголошує Ленін. Щоб перемогти, революційний клас у сфері свідомості повинен вимірювати себе не чужими мірками, а своїми, що визначаються його історичною місією. Відірвавши у буржуазії панування в цій сфері, озброївшись передовими досягненнями власної класової свідомості, пролетаріат стає дієвою матеріальною силою історичних перетворень: «Зброя критики, звісно, не може замінити критики зброєю, матеріальну силу потрібно повалити матеріальною ж силою, але теорія стає матеріальною силою, як тільки вона опановує маси», – писав Маркс у «До критики гегелівської філософії права».

Однак, досі ми говорили про суспільну свідомість не в умовах комунізму, а про стан речей, від якого потрібно відштовхуватися під час переходу до нього. На чому ж наполягає комунізм як наука щодо суспільної свідомості?

Комунізм наполягає на подоланні суспільного поділу праці, відчуження людини від інших людей, від процесу та результатів її діяльності. А це, у свою чергу, передбачає відповідні зміни в галузі суспільної свідомості. Відчуження форм суспільної свідомості як між собою, так і безпосередньо від людської діяльності, також буде подолано.

Ідеальні творчі здібності людини, які «замкнуті» в рамках форм суспільної свідомості, як окремих від матеріального виробництва сферах естетичної, моральної та інтелектуальної діяльності, стануть надбанням усіх і кожного.

Питання про суспільну свідомість за комунізму – це питання практики. Це питання планування та реалізації запланованого в матеріальній діяльності таким чином, щоб відбувалися відповідні зміни і в духовній сфері. Змінюючи свою діяльність, люди змінюють і свою свідомість як ідеальну сторону цієї діяльності. Але сама зміна практики за комунізму має свідомий характер. Тому тут вже не просто суспільне буття визначає суспільну свідомість, а суспільна свідомість, як свідомість злагоджено діючої вільної асоціації людей, визначає буття, яке визначає цю свідомість.

Усуспільнення всієї духовної культури, причому відповідно до логіки цієї культури, стає нагальною необхідністю для переходу до комунізму. Тому і в теорії комунізму проблеми суспільної свідомості мають сенс лише за умови, якщо їх розуміти як питання становлення комунізму.

У комуністичній діяльності, на відміну від відчуженої праці, ідеальне, як таке, що відповідає ідеалу, та ідеальне, як таке, що належить до «світу ідей», зрештою (саме зрештою) емпірично збігаються. Таким чином, ідеальне виступає як відображення (уявлення, вираження) матеріального в матеріальному – діяльне відображення об’єктивної реальності, і як ідеальна мета діяльності, з якою вона має узгоджуватися, і як естетичний ідеал. У цьому полягає найважливіший момент подолання роз’єднання праці, її відчуженості у всіх сенсах цього слова, і підтверджується самоцінність діяльності незалежно від того, що саме робиться.

Тут варто зазначити, що було б неправильно стверджувати, що свідомий і матеріальний аспекти людської діяльності є рівноправними в самій діяльності, незалежно від її історичної специфіки, і тим більше – на рівні діяльності суспільства як цілого. Звісно, ідеальний аспект присутній у будь-якій діяльності як аспект цілепокладання, якщо тільки цілепокладання не лежить поза межами цієї діяльності (для працівника, який лише продає свою робочу силу, його власна мета зовсім не збігається з метою виробництва), але йдеться про інше. У будь-якому випадку, спосіб виробництва матеріального життя визначає спосіб виробництва ідеального. І з цим доводиться рахуватися.

Це важливо з практичної точки зору, тому що, хоча ідея й стає матеріальною силою, коли опановує маси, сама вона є результатом усвідомлення матеріальних сил у процесі їхньої дії (результат і сторона взаємодії матеріального з матеріальним). Іншими словами, щоб змінити спосіб виробництва ідей, потрібно змінити спосіб виробництва матеріального життя. Лише так, а не навпаки. Навпаки нічого не вийде.

Наприклад, скільки б ми зараз не пропагували (тобто, діяли б ідеально) діалектику, від цього вона не стане звичним способом мислення, а виробництво діалектичного мислення й надалі залишатиметься, якщо можна так висловитися, ремеслом. Це ремесло, звісно, дуже важливе, і від нього значною мірою залежить майбутнє. Але ми повинні розуміти, що люди, які мислять діалектично, не стануть масовим продуктом, якщо збережеться нинішній спосіб виробництва матеріального життя і відповідний йому спосіб виробництва свідомості.

Щодо проблеми комунізму, необхідно не просто констатувати напрямок переходу матеріального життя людини в свідомість (що первинне, а що вторинне), але й запропонувати ідеї щодо практичного здійснення цього переходу. І тут надзвичайно важливі конкретні ідеї з цього приводу. Важливі ідеї щодо формування тих чи інших емпіричних форм комунізму, оскільки ніколи все матеріальне не переходить у свідомість миттєво, а завжди робить це в певних, особливих формах. Важливі, наприклад, такі ідеї, як у Макаренка: любов – це звичайна річ, яку потрібно правильно організувати.

Яку б величезну роль не відігравав ідеальний момент діяльності, йдеться завжди про життя та діяльність матеріальної людини в матеріальному світі. Ідеальне завжди є атрибутом матеріальної діяльності суспільної людини і відповідає її розвитку. «Царство свободи», тотальний комунізм, характеризується, зокрема, і свідомим виробництвом суспільної свідомості в суспільному масштабі.

Окремі прориви у сфері суспільної свідомості, які виходять за межі капіталізму, відбуваються вже зараз. Про них (зокрема) йшлося у другій главі цієї роботи. Це виробництво ідеального під час створення сучасних технологій. Але справа не лише в цьому, не кожне ідеальне належить до суспільної свідомості. Це вже не просто ідеальна діяльність у своїй відчуженості, яка має надбудовний характер, і не замкнуте в рамках професійного кретинізму ідеальне. Це ідеальний момент виробництва людського життя взагалі. І те, як саме це життя виробляється, визначає те, (як і у випадку з базами даних та вільним програмним забезпеченням, відбувається один і той самий процес) як створюється нове безпосереднє, невідчужене суспільне усвідомлення, хоч поки що на локальних ділянках. Це таке усвідомлення, яке формується не лише і не стільки самим фактом створення програми, скільки самим характером її створення та функціонування.

Суспільна свідомість може визначати суспільну свідомість лише через визначення суспільної практики. Причому останню не можна розуміти як щось відокремлене від цілісної людини, як щось ідеальне взагалі, а лише як момент діяльності матеріальної людини в матеріальному світі. Річ у тім, що справа полягає саме у свідомому виробництві цілісної людини – універсально діючої, чутливої, мислячої. Така діяльність, незалежно від того, що саме робиться, виступатиме як здійснення комунізму в безпосередньо колективній формі і буде діяльністю з виробництва свободи, цілепокладання, чуттєвості, моральності та мислення.

В основному, це передбачає, що будь-яка діяльність за комунізму повинна, зокрема, бути виробництвом ідей та ідеалу майбутнього життя, виробництвом людини.
\end{document}
